
\subsubsection{6.1 Interpreting the 53\% Scale Confound}

The primary finding --- that MMLU scale control reduces the TruthfulQA $\times$ bias-proxy correlation by 53.1\% --- does not imply that factuality and bias avoidance are unrelated. The residual partial correlation (0.343, p = 1.40 $\times 10^{-9})$ is positive and highly significant. The finding is that the raw leaderboard correlation (0.732) substantially overstates the alignment-specific component of the co-movement.

The confounding mechanism is straightforward: models with higher MMLU scores tend to be larger and trained on more data, performing better on factuality questions (TruthfulQA) and making fewer errors on reasoning-intensive bias-detection items (ARC proxy) simultaneously, not because their alignment training is coherent, but because general capability scales with both outcomes. MMLU as a common cause induces positive correlation between TruthfulQA and the bias proxy that is unrelated to alignment-specific training effects.

A practical implication for leaderboard interpretation: a model family that improves both TruthfulQA and BBQ-type scores may simply be scaling general capability. Attributing correlated cross-model improvement to aligned training objectives requires, at minimum, controlling for a general capability proxy such as MMLU.

\subsubsection{6.2 The Residual Partial Correlation}

After MMLU control, partial\_rho = 0.343 [0.180, 0.492] remains strongly significant. Three competing explanations are relevant:

\textbf{(1) ARC proxy artifact.} ARC Challenge measures logical reasoning; TruthfulQA also has a reasoning component requiring models to assess the truth value of plausible-sounding statements. The partial correlation 0.343 may largely reflect shared reasoning demands between these two benchmarks rather than any genuine factuality-bias alignment coupling. This explanation is consistent with the null RLHF sign test result (Tier 3, p = 0.686): if genuine alignment co-training drove the residual, one would expect RLHF models to consistently outperform their base counterparts on the bias proxy, which is not observed.

\textbf{(2) Genuine alignment co-training.} If RLHF and safety-focused fine-tuning jointly improve factuality and bias avoidance, cross-model variation in training intensity could generate residual coupling. The Tier 3 null result does not support this under the ARC proxy, though the proxy may be inadequate to capture RLHF-targeted bias behavior.

\textbf{(3) Residual family-level clustering.} Within-family training recipe variation not captured by the family-weighted correction could maintain residual positive correlation.

Disambiguation requires genuine HELM Lite BBQ per-model accuracy. If partial\_rho under real BBQ data is substantially lower than 0.343, explanation (1) is dominant. Replication with genuine BBQ data is the most important next step.

\subsubsection{6.3 Limitations}

\textbf{L1: BBQ proxy (ARC Challenge).} The substitution of ARC Challenge accuracy for HELM Lite BBQ per-model scores is the most consequential limitation. The primary finding --- that MMLU confounds alignment benchmark correlations (Fisher z p = 3.22 $\times 10^{-12})$ --- is methodology-valid regardless of proxy choice, as partial Spearman correctly removes scale variance from both benchmarks before computing rank correlation. Secondary claims about ''factuality-bias coupling'' are not interpretable as such and must be qualified as reflecting TruthfulQA $\times$ ARC correlation after MMLU control. The match\_rate = 1.000 is also an artifact of the same-source proxy, which means the cross-source join mechanics are not confirmed for a genuinely independent bias benchmark.

\textbf{L2: Scenario classification underpowered (N=296).} partial\_rho = 0.343 with CI [0.180, 0.492] spans the 0.40 structural threshold. Resolving the scenario assignment requires N $\geq$ 400--600 (per power calculations for distinguishing rho = 0.35 from rho = 0.41 at $\alpha$ = 0.05) or genuine BBQ data with different noise characteristics. The ambiguous outcome is pre-registered and scientifically valid but forecloses definitive structural conclusions.

\textbf{L3: Safety dimension excluded (HarmBench N=0).} The three-benchmark partial Spearman matrix (factuality $\times$ bias $\times$ safety) could not be computed. Model name incompatibility between LLM Leaderboard v1 and HarmBench Table 2 is the likely cause. The study is limited to two alignment dimensions.

\textbf{L4: RLHF sign test on proxy.} The null result (p = 0.686) may reflect proxy inadequacy: RLHF alignment targets human-perceived helpfulness, harmlessness, and honesty --- properties that include bias avoidance --- but ARC Challenge may not be sensitive to this RLHF signal.

\textbf{L5: Observational cross-sectional design.} The study characterizes population-level correlation structure for open-weight LLMs from the LLM Leaderboard v1 era (approximately 2022--2024). Findings do not generalize to proprietary models, post-2024 training regimes, or within-model training dynamics.

\subsubsection{6.4 Implications for Alignment Evaluation Practice}

The Fisher z difference test between raw and partial Spearman should be applied as a standard diagnostic step in alignment benchmark analysis whenever a capability proxy co-varies with the benchmarks of interest. The test is computationally inexpensive, requires only a single scale covariate, and provides a direct quantification of the confound magnitude. Positive correlations between alignment benchmarks should not be interpreted as evidence of alignment-specific coherence without first verifying that scale confounding is not the dominant driver.

